\phantomsection
\addcontentsline{toc}{section}{About Part IA}
\section*{About Part IA}

The first year is the only part of the Tripos with no choices in it. Everybody sits the
same eight courses, and the reason is that the year has a job to do: it has to take a
room of people who arrived knowing how to differentiate and turn them into people who
can write a proof. Almost everything in this part is either a foundation that later
years will build on without comment, or a piece of technique that later years will
assume you can execute in your sleep.

The year splits fairly cleanly in two. Numbers and Sets, Groups, Analysis I and Vectors
and Matrices are the pure half, and between them they introduce essentially every kind
of argument the rest of the book uses. Numbers and Sets comes first in more than just
the ordering -- it is where induction, the Euclidean algorithm, equivalence relations,
countability and the least upper bound axiom appear, and it is the course whose real
content is the idea of proof itself. Analysis I then takes the least upper bound axiom
and does something startling with it, building limits, continuity, differentiation and
the Riemann integral out of nothing but $\epsilon$ and $N$. Groups and Vectors and
Matrices run in parallel as the first taste of abstraction: Groups is abstract from its
first line and stays concrete in its examples, working through the dihedral, symmetric
and M\"obius groups by hand, while Vectors and Matrices spends most of its time on
$\R^3$ and $\C^n$ before quietly showing you that a matrix is a linear map in
disguise.

The applied half -- Differential Equations, Vector Calculus, Dynamics and Relativity,
and Probability -- is deliberately less careful, and it is worth being clear that this
is a choice rather than an oversight. Differential Equations solves equations without
ever proving that solutions exist. Vector Calculus differentiates under integral signs
and exchanges limits with abandon. Probability defines a random variable as a function
on a sample space and then never worries about whether the sample space is big enough.
All three of these debts are called in later, in Analysis and Topology, in Methods, and
in Probability and Measure respectively; the first year buys fluency now and pays for
rigour later.

There are through-lines worth watching for. Vector Calculus is the most quietly
important course of the year, because the divergence theorem and Stokes' theorem are
the entire grammar of Electromagnetism and Fluid Dynamics, and its chapter on curves
and surfaces is the first draft of everything in Geometry and Differential Geometry.
Groups is the seed of the longest chain in the book, running through Groups, Rings and
Modules to Galois Theory and Representation Theory. Vectors and Matrices is the same
material as Linear Algebra with the abstraction removed, which makes reading the two in
succession an unusually good way to see what abstraction actually buys you. And
Dynamics and Relativity, which looks at first like a physics course that wandered into
the wrong degree, turns out to be where the variational and Hamiltonian ideas of the
next two years get their motivation.

If you are reading this book rather than sitting the exams, Part IA is the one part
that genuinely rewards being read in order.
